% ЗГЕНЕРОВАНО scripts/build_format_v2_chapter.py — не редагувати вручну.
% Джерело: editorial/cards/11.*.json, format-v2/rozdil-11.json, paraleli.json.

\naukchapter{11}{8}{Глава 11. Переконання}
\naukintro{Одинадцята глава має одну будову на всі вісім віршів: цитата переконання — і наслідок для того, хто його має. niścayī, той, хто має переконання, повторюється на початку кожного вірша, і саме цей повтор і є формою глави. Тут важливо, що переконання не треба доводити: текст описує не аргумент, а стан того, хто вже так думає.}
\vsec{11.1}
\vdeva{भावाभावविकारश् च स्वभावाद् इति निश्चयी । निर्विकारो गतक्लेशः सुखेनैवोपशाम्यति ॥}
\viast{bhāvābhāvavikāraś ca svabhāvād iti niścayī | nirvikāro gatakleśaḥ sukhenaivopaśāmyati ||}
\vtranslit{бгавабгававікараш ча свабгавад іті нішчаї | нірвікаро ґатаклешах сукхенайвопашам'яті ||}
\vperek{«Видозміна буття й небуття — за власною природою» — хто має таке переконання, той без видозмін, з минулими турбаціями, втихає легко.}
\vcom{Построкове зіставлення: редакторська чернетка, не рецензовано.}
\vcom{Рядок а}
\vdeva{भावाभावविकारश् च स्वभावाद् इति निश्चयी}
\viast{bhāvābhāvavikāraś ca svabhāvād iti niścayī}
\vtranslit{бгавабгававікараш ча свабгавад іті нішчаї}
\vperek{«Видозміна буття й небуття — за власною природою» — хто має таке переконання,}
\vcom{bhāva й abhāva повʼязані з vikāra; внутрішнє членування складного слова потребує фахової звірки. svabhāvād — «за власною природою», iti niścayī — «той, хто має таке переконання».}
\vcom{Рядок б}
\vdeva{निर्विकारो गतक्लेशः सुखेनैवोपशाम्यति}
\viast{nirvikāro gatakleśaḥ sukhenaivopaśāmyati}
\vtranslit{нірвікаро ґатаклешах сукхенайвопашам'яті}
\vperek{той без видозмін, з минулими турбаціями, втихає легко.}
\vcom{nirvikāraḥ означає відсутність власної видозміни; gata-kleśaḥ — згаслі турбації; sukhena eva upaśāmyati — втишення саме з легкістю. Наслідок переконання в другому рядку граматично продовжує niścayī з першого.}
\vcom{`niścayī` («той, хто має переконання») повторюється на початку 11.1–11.4 — це структурна формула всієї глави, і переклад має тримати її однаково. `nirvikāraḥ` уточнює: не «без змін» у сенсі нерухомості, а без **власних** видозмін від того, що змінюється назовні. `sukhena eva` («саме з легкістю») — втишення не коштує зусиль, і це важливо.}
\vpar{Українська паралель (наша інтерпретація)}{Григорій Сковорода, «Потоп зміїний» (трактат).}{Наша паралель: у Сковороди видозміна речей теж відбувається за власною природою, і саме тому переконання про це дає спокій. Наше порівняння — типологічне, не твердження про джерело.}
\vfig{ilustratsii/11-1-motion.pdf}{Буття й небуття як видозміна за власною природою}{Схема наша, не історична ілюстрація.}

\vsec{11.2}
\vdeva{ईश्वरः सर्वनिर्माता नेहान्य इति निश्चयी । अन्तर्गलितसर्वाशः शान्तः क्वापि न सज्जते ॥}
\viast{īśvaraḥ sarvanirmātā nehānya iti niścayī | antargalitasarvāśaḥ śāntaḥ kvāpi na sajjate ||}
\vtranslit{ішварах сарванірмата неханя іті нішчаї | антарґалітасарвашах шантах квапі на саджджате ||}
\vperek{«Творець усього — Ішвара, і немає тут іншого» — хто має таке переконання, той із розчиненими всередині сподіваннями, втихлий, ніде не застрягає.}
\vcom{Построкове зіставлення: редакторська чернетка, не рецензовано.}
\vcom{Рядок а}
\vdeva{ईश्वरः सर्वनिर्माता नेहान्य इति निश्चयी}
\viast{īśvaraḥ sarvanirmātā nehānya iti niścayī}
\vtranslit{ішварах сарванірмата неханя іті нішчаї}
\vperek{«Творець усього — Ішвара, і немає тут іншого» — хто має таке переконання,}
\vcom{īśvaraḥ («Ішвара»), sarvanirmātā («творець усього») та nehānya («тут немає іншого») формують перше твердження. iti niścayī означає «той, хто має таке переконання» — дослівно повторювану граматичну рамку 11.1–11.8. Рядок описує позицію мовця, а не доводить конкретне історичне богослов'я.}
\vcom{Рядок б}
\vdeva{अन्तर्गलितसर्वाशः शान्तः क्वापि न सज्जते}
\viast{antargalitasarvāśaḥ śāntaḥ kvāpi na sajjate}
\vtranslit{антарґалітасарвашах шантах квапі на саджджате}
\vperek{той із розчиненими всередині сподіваннями, втихлий, ніде не застрягає.}
\vcom{antargalitasarvāśaḥ — складний прикметник про внутрішньо згаслі сподівання, śāntaḥ — «втихлий», kvāpi na sajjate — «ніде не чіпляється / не застрягає». Не ототожнювати образ внутрішньої свободи з байдужістю до обов'язків інших людей.}
\vcom{У 11.2 Ішвару (`īśvaraḥ`) прямо названо `sarvanirmātā` («творець усього») й заперечено іншого (`na iha anyaḥ`); сама назва Ішвари вже трапляється в 4.6 як `jagad īśvaram` (GRETIL) / `jagadīśvaram` (SARIT). Це не теологія, а спосіб зняти питання причини. `sajjate` — «застрягає, чіпляється»: наслідок не в тому, що людина стала всемогутньою, а в тому, що вона більше ні за що не тримається.}
\vpar{Українська паралель (наша інтерпретація)}{Митрополит Іларіон Київський, «Слово про закон і благодать» (XI ст.).}{Наша паралель: у київському Слові теж є думка, що за всім стоїть один початок, і тому сподівання людини не мають де застрягти. Наше порівняння — типологічне; датування цієї книги науково не встановлене.}
\vfig{ilustratsii/11-2-witness.pdf}{Один творець у центрі, і немає тут іншого}{Схема наша, не історична ілюстрація.}

\vsec{11.3}
\vdeva{आपदः संपदः काले दैवादेवेति निश्चयी । तृप्तः स्वस्थेन्द्रियो नित्यं न वाञ्छति न शोचति ॥}
\viast{āpadaḥ saṃpadaḥ kāle daivādeveti niścayī | tṛptaḥ svasthendriyo nityaṃ na vāñchati na śocati ||}
\vtranslit{ападах сампадах кале дайвадеветі нішчаї | тріптах свастхендрійо нітям на ваньччаті на шочаті ||}
\vperek{«Нещастя й достаток у свій час — саме від долі» — хто має таке переконання, той вдоволений, із чуттями в порядку, повсякчас не бажає й не сумує.}
\vcom{Построкове зіставлення: редакторська чернетка, не рецензовано.}
\vcom{Рядок а}
\vdeva{आपदः संपदः काले दैवादेवेति निश्चयी}
\viast{āpadaḥ saṃpadaḥ kāle daivādeveti niścayī}
\vtranslit{ападах сампадах кале дайвадеветі нішчаї}
\vperek{«Нещастя й достаток у свій час — саме від долі» — хто має таке переконання,}
\vcom{āpadaḥ та saṃpadaḥ позначають нещастя і достаток; kāle — «у свій час», daivād eva — «саме від долі». iti niścayī позначає переконання, як у 11.1–11.2. Це філософська позиція персонажа, не емпіричне твердження, що всі життєві події спричиняє доля.}
\vcom{Рядок б}
\vdeva{तृप्तः स्वस्थेन्द्रियो नित्यं न वाञ्छति न शोचति}
\viast{tṛptaḥ svasthendriyo nityaṃ na vāñchati na śocati}
\vtranslit{тріптах свастхендрійо нітям на ваньччаті на шочаті}
\vperek{той вдоволений, із чуттями в порядку, повсякчас не бажає й не сумує.}
\vcom{tṛptaḥ — «вдоволений», svasthendriyaḥ — «із чуттями у стані рівноваги», nityam — повсякчас. Пара na vāñchati / na śocati заперечує і прагнення, і сумування; український повтор «не ... й не ...» збережено без нової обіцянки абсолютної нечутливості.}
\vcom{`daiva` — «доля» як те, що розподіляє, а не як примха. Формула «не бажає й не сумує» — та сама пара, що в 8.1–8.2, і переклад має її впізнавати. `svasthendriya` («із чуттями в порядку») — рідкісний складник: ідеться не про приборкання чуттів, а про їхню справність.}
\vpar{Українська паралель (наша інтерпретація)}{Пантелеймон Куліш, «Чорна рада» (1857).}{Наша паралель: у Куліша доля теж постає як те, що роздає нещастя й достаток не за заслугою. Наше порівняння — за будовою думки.}
\vfig{ilustratsii/11-3-balance.pdf}{Дві шальки: нещастя й достаток у свій час}{Схема наша, не історична ілюстрація.}

\vsec{11.4}
\vdeva{सुखदुःखे जन्ममृत्यू दैवादेवेति निश्चयी । साध्यादर्शी निरायासः कुर्वन्नपि न लिप्यते ॥}
\viast{sukhaduḥkhe janmamṛtyū daivādeveti niścayī | sādhyādarśī nirāyāsaḥ kurvannapi na lipyate ||}
\vtranslit{сукхадухкхе джанмамрітю дайвадеветі нішчаї | садгядарші ніраясах курваннапі на ліп'яте ||}
\vperek{«Щастя й страждання, народження й смерть — саме від долі» — хто має таке переконання, той, не бачачи нічого досяжного, без натуги, навіть діючи, не забруднюється.}
\vcom{Построкове зіставлення: редакторська чернетка, не рецензовано.}
\vcom{Рядок а}
\vdeva{सुखदुःखे जन्ममृत्यू दैवादेवेति निश्चयी}
\viast{sukhaduḥkhe janmamṛtyū daivādeveti niścayī}
\vtranslit{сукхадухкхе джанмамрітю дайвадеветі нішчаї}
\vperek{«Щастя й страждання, народження й смерть — саме від долі» — хто має таке переконання,}
\vcom{sukhaduḥkhe — двоїна «щастя і страждання», janmamṛtyū — пара «народження і смерть», daivād eva — «саме від долі». Рефрен iti niścayī повторює формулу 11.3; перелік доповнює попереднє зіставлення протилежностей.}
\vcom{Рядок б}
\vdeva{साध्यादर्शी निरायासः कुर्वन्नपि न लिप्यते}
\viast{sādhyādarśī nirāyāsaḥ kurvannapi na lipyate}
\vtranslit{садгядарші ніраясах курваннапі на ліп'яте}
\vperek{той, не бачачи нічого досяжного, без натуги, навіть діючи, не забруднюється.}
\vcom{sādhyādarśī потребує обережної передачі як «той, хто не бачить досяжної мети»; nirāyāsaḥ — «без натуги», kurvann api na lipyate — «навіть діючи, не забруднюється». Тут розрізнено відсутність чіпляння за результат і фізичну відмову від будь-якої дії.}
\vcom{`sādhyādarśī` («той, хто не бачить досяжного») — найтонше слово глави: дія не припиняється, але зникає те, заради чого вона нібито робиться. `kurvann api na lipyate` — «навіть діючи, не забруднюється»: це не відмова від дії, а її розпредметнення. Перелік розширюється від 11.3 до 11.4 (додано народження й смерть), і це наростання навмисне.}
\vpar{Українська паралель (наша інтерпретація)}{Тарас Шевченко, «Давидові псалми» (1845).}{Наша паралель: у Шевченкових переспівах псалмів щастя й страждання теж не міряються людською міркою. Наше порівняння — типологічне.}
\vfig{ilustratsii/11-4-balance.pdf}{Дві шальки: народження й смерть, щастя й страждання}{Схема наша, не історична ілюстрація.}

\vsec{11.5}
\vdeva{चिन्तया जायते दुःखं नान्यथेहेति निश्चयी । तया हीनः सुखी शान्तः सर्वत्र गलितस्पृहः ॥}
\viast{cintayā jāyate duḥkhaṃ nānyatheheti niścayī | tayā hīnaḥ sukhī śāntaḥ sarvatra galitaspṛhaḥ ||}
\vtranslit{чінтая джаяте духкхам нанятхехеті нішчаї | тая хінах сукхі шантах сарватра ґалітаспріхах ||}
\vperek{«Страждання народжується з думання, не інакше» — хто має таке переконання, той, позбавлений його, щасливий, втихлий, скрізь із опалою спрагою.}
\vcom{Построкове зіставлення: редакторська чернетка, не рецензовано.}
\vcom{Рядок а}
\vdeva{चिन्तया जायते दुःखं नान्यथेहेति निश्चयी}
\viast{cintayā jāyate duḥkhaṃ nānyatheheti niścayī}
\vtranslit{чінтая джаяте духкхам нанятхехеті нішчаї}
\vperek{«Страждання народжується з думання, не інакше» — хто має таке переконання,}
\vcom{cintayā — орудний від cintā: турботливе мисленнєве обертання; jāyate duḥkhaṃ — «народжується страждання», nānyathā — «не інакше». Рефрен iti niścayī подає це як висловлене переконання, а не універсальне медичне пояснення страждання.}
\vcom{Рядок б}
\vdeva{तया हीनः सुखी शान्तः सर्वत्र गलितस्पृहः}
\viast{tayā hīnaḥ sukhī śāntaḥ sarvatra galitaspṛhaḥ}
\vtranslit{тая хінах сукхі шантах сарватра ґалітаспріхах}
\vperek{той, позбавлений його, щасливий, втихлий, скрізь із опалою спрагою.}
\vcom{tayā hīnaḥ — «позбавлений її» (попередньої cintā), sukhī śāntaḥ — «щасливий, втихлий», sarvatra galitaspṛhaḥ — «скрізь із згаслим прагненням». Це не твердження про шкідливість інтелектуальної праці; в контексті йдеться про прив'язливе хвилювання.}
\vcom{`cintā` — не «думка» взагалі, а **турботливе обертання думки**, і саме воно назване джерелом страждання. Переклад «думання» тут кращий за «думка», бо остання звучить нейтрально. `galitaspṛha` («із опалою спрагою») — образ, який текст уже вживав у 11.2 (`antargalita`); словник лишається спільним.}
\vpar{Українська паралель (наша інтерпретація)}{Памфіл Юркевич, «Серце і його значення в духовному житті людини» (1860).}{Наша паралель: коли вірш каже, що страждання народжується з думання, він ставить джерело болю в думку, а не в річ. Юркевич шукає це джерело в серці — наше порівняння за способом ставити питання.}
\vfig{ilustratsii/11-5-flame.pdf}{Полумʼя, що живиться думанням}{Схема наша, не історична ілюстрація.}

\vsec{11.6}
\vdeva{नाहं देहो न मे देहो बोधो ऽहम् इति निश्चयी । कैवल्यम् इव संप्राप्तो न स्मरत्य् अकृतं कृतम् ॥}
\viast{nāhaṃ deho na me deho bodho 'ham iti niścayī | kaivalyam iva saṃprāpto na smaraty akṛtaṃ kṛtam ||}
\vtranslit{нахам дехо на ме дехо бодго ’хам іті нішчаї | кайвалям іва сампрапто на смаратй акрітам крітам ||}
\vperek{«Я не тіло, і тіло не моє; я — усвідомлення» — хто має таке переконання, той, наче досягши окремішності, не пам'ятає незробленого як зробленого.}
\vcom{Построкове зіставлення: редакторська чернетка, не рецензовано.}
\vcom{Рядок а}
\vdeva{नाहं देहो न मे देहो बोधो ऽहम् इति निश्चयी}
\viast{nāhaṃ deho na me deho bodho 'ham iti niścayī}
\vtranslit{нахам дехо на ме дехо бодго ’хам іті нішчаї}
\vperek{«Я не тіло, і тіло не моє; я — усвідомлення» — хто має таке переконання,}
\vcom{nāhaṃ deho, na me deho — два заперечення ототожнення із тілом та володіння ним; bodho 'ham — «я — усвідомлення». iti niścayī повторює рамку глави; твердження стосується самопізнання в адвайті, а не фактичної відсутності тіла.}
\vcom{Рядок б}
\vdeva{कैवल्यम् इव संप्राप्तो न स्मरत्य् अकृतं कृतम्}
\viast{kaivalyam iva saṃprāpto na smaraty akṛtaṃ kṛtam}
\vtranslit{кайвалям іва сампрапто на смаратй акрітам крітам}
\vperek{той, наче досягши окремішності, не пам'ятає незробленого як зробленого.}
\vcom{kaivalyam iva saṃprāptaḥ буквально містить iva («наче»), тому стан не проголошено беззастережно досягнутим. na smaraty akṛtaṃ kṛtam — стислий і неоднозначний вислів про зроблене й незроблене; чинний український переклад дослівно збережено, але він потребує фахового зіставлення, а не вигаданої однозначної психологічної інтерпретації.}
\vcom{Формула «не тіло, і тіло не моє» та сама, що в 8.4 і 2.22 — тут вона подана як переконання, а не як висновок. `kaivalya` («окремішність») з'являється один раз і з `iva` («наче»): текст не каже, що стан досягнуто. `akṛtaṃ kṛtam` — «незроблене як зроблене»: ідеться не про забудькуватість, а про те, що минуле більше не рахується.}
\vpar{Українська паралель (наша інтерпретація)}{Григорій Сковорода, «Наркіс. Розмова про те: Пізнай себе».}{Наша паралель: у Сковороди справжня людина теж не збігається з тілом і не належить йому. Наше порівняння — типологічне, не про вплив.}
\vfig{ilustratsii/11-6-neti.pdf}{Тіло й володіння тілом відкинуті: лишається усвідомлення}{Схема наша, не історична ілюстрація.}

\vsec{11.7}
\vdeva{आब्रह्मस्तम्बपर्यन्तम् अहम् एवेति निश्चयी । निर्विकल्पः शुचिः शान्तः प्राप्ताप्राप्तविनिर्वृतः ॥}
\viast{ābrahmastambaparyantam aham eveti niścayī | nirvikalpaḥ śuciḥ śāntaḥ prāptāprāptavinirvṛtaḥ ||}
\vtranslit{абрахмастамбапар'янтам ахам еветі нішчаї | нірвікалпах шучіх шантах праптапраптавінірврітах ||}
\vperek{«Від Брахми до стеблини — це саме я» — хто має таке переконання, той без домислювання, чистий, втихлий, вдоволений і досягнутим, і недосягнутим.}
\vcom{Построкове зіставлення: редакторська чернетка, не рецензовано.}
\vcom{Рядок а}
\vdeva{आब्रह्मस्तम्बपर्यन्तम् अहम् एवेति निश्चयी}
\viast{ābrahmastambaparyantam aham eveti niścayī}
\vtranslit{абрахмастамбапар'янтам ахам еветі нішчаї}
\vperek{«Від Брахми до стеблини — це саме я» — хто має таке переконання,}
\vcom{ābrahmastambaparyantam окреслює межу «від Брахми до стеблини», aham eva — «це саме я», iti niścayī — «хто має таке переконання». Це метафізична формула всеохопності, а не біологічне ототожнення окремої людини з рослиною.}
\vcom{Рядок б}
\vdeva{निर्विकल्पः शुचिः शान्तः प्राप्ताप्राप्तविनिर्वृतः}
\viast{nirvikalpaḥ śuciḥ śāntaḥ prāptāprāptavinirvṛtaḥ}
\vtranslit{нірвікалпах шучіх шантах праптапраптавінірврітах}
\vperek{той без домислювання, чистий, втихлий, вдоволений і досягнутим, і недосягнутим.}
\vcom{nirvikalpaḥ — «без домислювання або розрізнювальних конструкцій», śuciḥ і śāntaḥ — «чистий і втихлий»; prāptāprāptavinirvṛtaḥ об'єднує «досягнуте» й «недосягнуте». Чинну українську пару збережено як етичний образ неприв'язаності до набуття.}
\vcom{`ābrahmastambaparyantam` — та сама формула межі, що в 2.11 (`brahmādistambaparyantam`), і переклад має бути однаковий. `prāptāprāpta-vinirvṛta` — «вдоволений і досягнутим, і недосягнутим»: це не байдужість, а відсутність різниці між двома станами. `śuci` тут — не ритуальна чистота, а відсутність домішки, бо стоїть поруч із `nirvikalpa`.}
\vpar{Українська паралель (наша інтерпретація)}{Іван Вишенський, «Книжка» (кінець XVI ст.).}{Наша паралель: у Вишенського межа між високим і низьким у світі теж знімається перед тим, що одне. Наше порівняння — типологічне.}
\vfig{ilustratsii/11-7-neti.pdf}{Від Брахми до стеблини — це саме я}{Схема наша, не історична ілюстрація.}

\vsec{11.8}
\vdeva{नानाश्चर्यम् इदं विश्वं न किञ्चिद् इति निश्चयी । निर्वासनः स्फूर्तिमात्रो न किञ्चिद् इव शाम्यति ॥}
\viast{nānāścaryam idaṃ viśvaṃ na kiñcid iti niścayī | nirvāsanaḥ sphūrtimātro na kiñcid iva śāmyati ||}
\vtranslit{нанашчар'ям ідам вішвам на кіньчід іті нішчаї | нірвасанах спхуртіматро на кіньчід іва шам'яті ||}
\vperek{«Цей світ, різноманітно дивовижний, — аніщо» — хто має таке переконання, той без усталених схильностей, саме лише сяйнування, втихає, наче аніщо.}
\vcom{Построкове зіставлення: редакторська чернетка, не рецензовано.}
\vcom{Рядок а}
\vdeva{नानाश्चर्यम् इदं विश्वं न किञ्चिद् इति निश्चयी}
\viast{nānāścaryam idaṃ viśvaṃ na kiñcid iti niścayī}
\vtranslit{нанашчар'ям ідам вішвам на кіньчід іті нішчаї}
\vperek{«Цей світ, різноманітно дивовижний, — аніщо» — хто має таке переконання,}
\vcom{nānāścaryam idaṃ viśvaṃ — «цей різноманітно дивовижний світ», na kiñcid — «аніщо»; рамка iti niścayī завершує восьмикратний повтор глави 11. Українське «аніщо» є перекладацьким переданням філософського судження, не запереченням емпіричних об'єктів.}
\vcom{Рядок б}
\vdeva{निर्वासनः स्फूर्तिमात्रो न किञ्चिद् इव शाम्यति}
\viast{nirvāsanaḥ sphūrtimātro na kiñcid iva śāmyati}
\vtranslit{нірвасанах спхуртіматро на кіньчід іва шам'яті}
\vperek{той без усталених схильностей, саме лише сяйнування, втихає, наче аніщо.}
\vcom{nirvāsanaḥ — «без закріплених схильностей», sphūrtimātro — «саме лише сяйнування / виявлення», na kiñcid iva śāmyati — «втихає, наче аніщо». Частка iva («наче») має залишитися помітною: вірш не оголошує людину буквально неіснуючою.}
\vcom{Глава закінчується тим самим словом, що й почалася в 11.1: `niścayī`, але тепер воно стоїть при «аніщо». `nirvāsana` («без усталених схильностей») — те саме `vāsanā`, що в 9.8; словник спільний. `sphūrtimātra` («саме лише сяйнування») і `na kiñcid iva` («наче аніщо») разом дають обережність, яку легко втратити: текст не каже, що людина зникає, — він каже «наче».}
\vpar{Українська паралель (наша інтерпретація)}{Богдан-Ігор Антонич, «Зелена Євангелія» (1938).}{Наша паралель: в Антонича світ теж постає дивовижним і водночас не вартим того, щоб за нього триматися. Наше порівняння — за образом.}
\vfig{ilustratsii/11-8-illusion.pdf}{Дивовижний світ під серпанком: аніщо}{Схема наша, не історична ілюстрація.}

